%%% Hlavní soubor. Zde se definují základní parametry a odkazuje se na ostatní části. %%%

%% Verze pro jednostranný tisk:
% Okraje: levý 40mm, pravý 25mm, horní a dolní 25mm
% (ale pozor, LaTeX si sám přidává 1in)
\documentclass[12pt,a4paper]{report}

\setlength\textwidth{145mm}
\setlength\textheight{247mm}
\setlength\oddsidemargin{15mm}
\setlength\evensidemargin{15mm}
\setlength\topmargin{0mm}
\setlength\headsep{0mm}
\setlength\headheight{0mm}
% \openright zařídí, aby následující text začínal na pravé straně knihy
\let\openright=\clearpage

%% Pokud tiskneme oboustranně:
% \documentclass[12pt,a4paper,twoside,openright]{report}
% \setlength\textwidth{145mm}
% \setlength\textheight{247mm}
% \setlength\oddsidemargin{15mm}
% \setlength\evensidemargin{0mm}
% \setlength\topmargin{0mm}
% \setlength\headsep{0mm}
% \setlength\headheight{0mm}
% \let\openright=\cleardoublepage

%% Použité kódování znaků: obvykle latin2, cp1250 nebo utf8:
\usepackage[utf8]{inputenc}

%% Ostatní balíčky
\usepackage{graphicx}
\usepackage{amsthm}


%% nastaveni 0 jako odskoceni prvniho radku, mezery mezi odstavci
\usepackage{parskip}

%% nastavení linespace na 1,5
\usepackage{setspace}
\onehalfspacing

%% umožnit eps v pdftex
\usepackage{epstopdf}



%% Balíček hyperref, kterým jdou vyrábět klikací odkazy v PDF,
%% ale hlavně ho používáme k uložení metadat do PDF (včetně obsahu).
%% POZOR, nezapomeňte vyplnit jméno práce a autora.
\usepackage[pdftex,unicode,hidelinks]{hyperref}   % Musí být za všemi ostatními balíčky
\hypersetup{pdftitle=Querying Preprocessed Linked Data Resources}
\hypersetup{pdfauthor=Petr Vávro}

%%% Drobné úpravy stylu

% Tato makra přesvědčují mírně ošklivým trikem LaTeX, aby hlavičky kapitol
% sázel příčetněji a nevynechával nad nimi spoustu místa. Směle ignorujte.
\makeatletter
\def\@makechapterhead#1{
  {\parindent \z@ \raggedright \normalfont
   \Huge\bfseries \thechapter. #1
   \par\nobreak
   \vskip 20\p@
}}
\def\@makeschapterhead#1{
  {\parindent \z@ \raggedright \normalfont
   \Huge\bfseries #1
   \par\nobreak
   \vskip 20\p@
}}
\makeatother

% Toto makro definuje kapitolu, která není očíslovaná, ale je uvedena v obsahu.
\def\chapwithtoc#1{
\chapter*{#1}
\addcontentsline{toc}{chapter}{#1}
}

\begin{document}

% Trochu volnější nastavení dělení slov, než je default.
\lefthyphenmin=2
\righthyphenmin=2

%%% Titulní strana práce

\pagestyle{empty}
\begin{center}

\large

Charles University in Prague

\medskip

Faculty of Mathematics and Physics

\vfill

{\bf\Large MASTER THESIS}

\vfill

\centerline{\mbox{\includegraphics[width=60mm]{../img/logo.eps}}}

\vfill
\vspace{5mm}

{\LARGE Bc. Petr Vávro}

\vspace{15mm}

% Název práce přesně podle zadání
{\LARGE\bfseries Querying Preprocessed Linked Data Resources}

\vfill

% Název katedry nebo ústavu, kde byla práce oficiálně zadána
% (dle Organizační struktury MFF UK)
Department of Software Engineering

\vfill

\begin{tabular}{rl}

Supervisor of the master thesis: & Mgr. Martin Nečaský, Ph.D. \\
\noalign{\vspace{2mm}}
Study programme: & Softwarové systémy I2 \\
\noalign{\vspace{2mm}}
Specialization: & Softwarové inženýrství \\
\end{tabular}

\vfill

% Zde doplňte rok
Prague 2014

\end{center}

\newpage

%%% Následuje vevázaný list -- kopie podepsaného "Zadání diplomové práce".
%%% Toto zadání NENÍ součástí elektronické verze práce, nescanovat.

%%% Na tomto místě mohou být napsána případná poděkování (vedoucímu práce,
%%% konzultantovi, tomu, kdo zapůjčil software, literaturu apod.)

\openright

\noindent
Dedication.

\newpage

%%% Strana s čestným prohlášením k diplomové práci

\vglue 0pt plus 1fill

\noindent
I declare that I carried out this master thesis independently, and only with the cited
sources, literature and other professional sources.

\medskip\noindent
I understand that my work relates to the rights and obligations under the Act No.
121/2000 Coll., the Copyright Act, as amended, in particular the fact that the Charles
University in Prague has the right to conclude a license agreement on the use of this
work as a school work pursuant to Section 60 paragraph 1 of the Copyright Act.

\vspace{10mm}

\hbox{\hbox to 0.5\hsize{%
In ........ date ............
\hss}\hbox to 0.5\hsize{%
signature of the author
\hss}}

\vspace{20mm}
\newpage

%%% Povinná informační strana diplomové práce

\vbox to 0.5\vsize{
\setlength\parindent{0mm}
\setlength\parskip{5mm}

Název práce:
Querying Preprocessed Linked Data Resources
% přesně dle zadání

Autor:
Petr Vávro

Katedra:  % Případně Ústav:
Katedra softwarového inženýrství
% dle Organizační struktury MFF UK

Vedoucí diplomové práce:
Mgr. Martin Nečaský, Ph.D.
% dle Organizační struktury MFF UK, případně plný název pracoviště mimo MFF UK

Abstrakt:
% abstrakt v rozsahu 80-200 slov; nejedná se však o opis zadání diplomové práce

Klíčová slova:
Linked Data, RDF, SPARQL, JSON, query optimization
% 3 až 5 klíčových slov

\vss}\nobreak\vbox to 0.49\vsize{
\setlength\parindent{0mm}
\setlength\parskip{5mm}

Title:
Querying Preprocessed Linked Data Resources
% přesný překlad názvu práce v angličtině

Author:
Petr Vávro

Department:
Department of Software Engineering
% dle Organizační struktury MFF UK v angličtině

Supervisor:
Mgr. Martin Nečaský, Ph.D.
% dle Organizační struktury MFF UK, případně plný název pracoviště
% mimo MFF UK v angličtině

Abstract:
% abstrakt v rozsahu 80-200 slov v angličtině; nejedná se však o překlad
% zadání diplomové práce

Keywords:
Linked Data, RDF, SPARQL, JSON, query optimization
% 3 až 5 klíčových slov v angličtině

\vss}

\newpage

%%% Strana s automaticky generovaným obsahem diplomové práce. U matematických
%%% prací je přípustné, aby seznam tabulek a zkratek, existují-li, byl umístěn
%%% na začátku práce, místo na jejím konci.

\openright
\pagestyle{plain}
\setcounter{page}{1}
\tableofcontents

%%% Jednotlivé kapitoly práce jsou pro přehlednost uloženy v samostatných souborech
\chapter*{Introduction}
\addcontentsline{toc}{chapter}{Introduction}


\chapter{Objectives}


\paragraph{Web service for storing Linked data resources}

Create a web service that ...

\paragraph{Searching of stored resources}
Searching

\paragraph{Ease of use}

\paragraph{Usability of the service}
Usability possibilities

\chapter{Thesis structure}
\label{chap:Structure}

\begin{itemize}
\item Chapter by chapter what is contained
\end {itemize}

\chapter{Analysis}
\label{chap:Analysis}
\todoi{add a chapter summary}

\section{Triplestore speed}
As briefly introduced in \autoref{chap:Introduction} triplestores are effective at processing queries. Their effectiveness originates in the specific indices they use \cite{onlineAllgeroGraphTripleIndices}. But in an application data are usually used in views and these get constructed from multiple subject-predicate-object triples, where more than one subject is present.
\todoi{add example of nested data structure}
As \cite{onlineDbisQueryOptimizationInRdf} states the main achievable optimization for a graph database data retrieval is through data locality - optimizing database pages so, that data frequently retrieved together is located near each other. But in a general triplestore, that provides data for multiple different views, data can't really be optimized this way, as the it is retrieved for queries that can and will overlap, but won't be identical in every aspect.
\todoi{example of anouther nested data structure that needs only part of the previous example}

That leads to the conclusion, that there have to be ways to retrieve the mentioned data views faster for an application.

\section{Storage requirements}
\label{sec:StorageRequirements}
The first way to speed up data retrieval that comes to mind is a query-results in-memory cache. Such a cache would store the result of the query along with the query and its parameters. On a subsequent query with the same parameters the cache would only retrieve the cached data and return it. This would skip all the steps needed to be executed by the triplestore - constructing an execution plan, executing the query and retrieving all the necessary database pages containing the requested data. This is a fine idea, but has its drawbacks.
%% maybe do as a list of items
\begin{itemize}
	\item First of it's a tradeoff of space for time. The same data stored in the query store can be retrieved by various queries and form different views. The cache would have to store all these different views separately. As the cache is in memory the amount of data will be constrained by its size. That can be partly alleviated by storing cache contents to a hard disk drive and compressing similar data.
	\item Secondly after the data changes the cache would still return old data. So we need a way to invalidate data in the cache, as on its own its detached from the original SPARQL data source and wouldn't know about the changes.
	\item Thirdly a simple cache doesn't know anything about the data view or the query that produced it and wouldn't provide anything else. If we add some meaning to the data and the query that produced it we could also build a searchable index over it.
\end{itemize}
But with a searchable index that is not a mere cache of query results anymore.

%% Data = resource -> resource has Id (how to get that id) -> first index B-Trees (already implemented better in Lucene)
For our needs the data resulted from the SPARQL query can be seen as a resource. To comply with the Linked Data principles this resource has a type and an id. It's type is defined by the query that produced it and it's id is defined by the parameters of the query. This provides us with a means to uniquely identify any cached data. At the same time it gives us the first two indexes we need to build. This defined resource has other beforehand unknown properties, any searchable property will give us another index.

\section{Resource Querying}
\label{sec:ResourceQuerying}
The basic way of getting the resource by its unique identifier was already described in \autoref{sec:StorageRequirements}. The resource has other properties witch could be searched as well. These properties are defined by the query that produced the data. In SPARQL there are two possible ways of creating a view of the data \cite{onlineSparqlQueryLanguageSpec}:
\begin{itemize}
	\item Describe Query - Used to extract an RDF graph from the SPARQL endpoint, the contents of which is left to the endpoint to decide based on what the maintainer deems as useful information.
	\item Construct Query - Used to extract information from the SPARQL endpoint and transform the results into valid RDF.
\end{itemize}
A describe query usually returns all the triples where the requested id is the subject or the object. As such it isn't really helpful in applications - as per specification it can return any data deemed useful.
Construct query specifies a transformation of the data. So for different parameters the results of the construct query should have same properties - return a similar graph of triplets. We expect that a graph, that represents a view of a resource and its properties, will contain a one root object and the other objects present in the graph are values of the root objects properties.
These properties can be extracted from the returned data and added to an index. A triple in RDF data as mentioned in \autoref{chap:Introduction} has three parts subject-predicate-object. For the resources as described in this thesis the predicate part equals to a property on the resource, the object is the value of the property. The basic recognized types in RDF for an object are \cite{onlineRdfConcepts}:
\begin{itemize}
	\item Literal - that consist of two or three elements: %%http://www.w3.org/TR/rdf11-concepts/#section-Graph-Literal
		\subitem Lexical form string
		\subitem Datatype IRI - being an IRI identifying a datatype that determines how the lexical form maps to a literal value
		\subitem Language tag - specified only when the the datatype IRI is a langString \footnote{'http://www.w3.org/1999/02/22-rdf-syntax-ns\#langString'}
	\item IRI
	\item blank node
\end{itemize}
Linked data promotes the use of vocabularies %%this is vague - be more precise
that further specify the objects type - specify the range of the known datatype IRI set.
\todoi{sample of a vocabulary entry}
Among the core datatypes of literals defined by RDF are string, number, date and timespan. These types enable querying our stored resources not only by their equality, but also gives the option to query by ranges. So the operators the created service should support depending on data type are:
\begin{itemize}
	\item String - Equals, Not Equals
	\item Number, Date, TimeSpan - Equals, Not Equals, Greater than, Less than, Between
%% didn't i forget some operators?
\end{itemize}
\todoi{in operator is an extension - maybe mention later - as multiple OR values are supported even in basic syntax}
Another useful operation for long texts is fulltext search. %% specify how a fulltext search usually works? - startswith, exact, suggestions etc.
This type of operation is at the current time not supported in SPARQL, and could be a major extension for usability of the service.

Search using basic operators could be implemented using a B-Tree stored on disk. As B-Trees support effective look up for specific data and range queries over the data.
%%here it sounds too easy - as databases have to also create a query plan - how to execute the query
Doing fulltext searching over data is a project on its own. Adding fulltext capabilities to applications is usually done using complex tools. In fulltext we usually not only require to find the data that contains the searched terms. Often fulltext searches have to do language specific analysis of the data and search terms. It is expected that the results of a fulltext query are returned in order of relevance.

Fulltext search and querying using the defined operators are supported in many database solutions. Specifically document databases use a document as their smallest entity. This document in a document database coincides with the defined term of a resource, which is the smallest entity we want to store and retrieve from the designed service.

\todoi{maybe add in introduction section about existing database solutions (graph db, document db, nosql vs sql)- where the term document database is introduced}
%% maybe use this somewhere?
%\section{Document databases}
%Opposite to the graph databases are usually considered to be document databases. Document databases consider the smallest entity to be the Domain aggregate which is stored as a whole document and always read at once. The biggest difference between these two database structures are the possibilities to query data
%\begin{itemize}
	%\item Graph databases provide ways to make complex queries over the whole data
	%\item Document databases usually provide only ways to query one document (data entity)
%\end{itemize}
%Document databases are nowadays considered to be the go to database for quick data retrieval.

\todoi{add document db comparison? - why choose raven}
%\section{Choosing a storage engine}
%- Document database comparisons
%\begin{itemize}
	%\item Document database options
%\end{itemize}


%%\begin{itemize}
%%	\item How to structure the data (JSON-LD)
%%		\subitem JSON-LD algorithms
%%		\subitem As whole document
%%			\subsubitem Flat
%%			\subsubitem Expanded
%%		\subitem Split to multiple documents
%%			\subsubitem How to split
%%\end{itemize}
%% isn't this more of a design choice than analysis?
\section{Resource storage options}
Available JSON serializations of RDF are \cite{onlineW3CRdfSyntax}:
\begin{itemize}
	\item JSON-LD
	\item RDF/JSON
	\item Rdfj
	\item jsonGRDDL
	\item JSON+RDF
\end{itemize}
\todoi{maybe add some discussion about the formats?}
As most of the available document databases use JSON as their storage format, it's desirable to select an JSON RDF serialization. And the most human readable and easiest to consume serialization format is JSON-LD, which was designed with these two goals in mind \cite{onlineJsonLdDesignGoals}. \todoi{source?}

\todoi{add json-ld serialization sample and some json-ld format description?}
There is more than one way to serialize RDF into JSON-LD. JSON-LD defines processing algorithms, that algorithms manipulate the look and size of the serialized data and try to create normalized versions of it \cite{onlineW3CJsonLdAlgorithmsSpec}:
\begin{description}
	\item[Expansion Algorithm] This algorithm expands a JSON-LD document, such that all context definitions are removed, all terms and compact IRIs are expanded to absolute IRIs, blank node identifiers, or keywords and all JSON-LD values are expressed in arrays in expanded form.
	\item[Compaction Algorithm] This algorithm compacts a JSON-LD document, such that the given context is applied. This must result in shortening any applicable IRIs to terms or compact IRIs, any applicable keywords to keyword aliases, and any applicable JSON-LD values expressed in expanded form to simple values such as strings or numbers.
	\item[Flattening Algorithm] This algorithm flattens an expanded JSON-LD document by collecting all properties of a node in a single JSON object and labeling all blank nodes with blank node identifiers. This resulting uniform shape of the document, may drastically simplify the code required to process JSON-LD data in certain applications.
\end{description}
\todoi{add samples for what the algorithms do - but its in the spec link}

When requesting data from a SPARQL endpoint that supports returning JSON-LD it's not possible to specify the algorithms the server should run on the returned data, so its up to the client to alter the serialization.

As stated in \autoref{sec:ResourceQuerying} the construct query, which will be used for getting the view of the resource, returns a graph of triplets/objects. This graph could be stored in multiple presentations:
\begin{description}
\item[Flat presentation] As result of the Flattening algorithm. The objects will all be listed in the root @graph property of the JSON-LD data serialization.
\item[Nested presentation] Where objects of the graph are nested into a root object and so form a hierarchy. This presentation follows a natural understanding of an object with properties that may have other complex objects as theirs values. This is not a normalization that is a result of applying JSON-LD algorithms.
\end{description}
Document databases, that use JSON as their data serialization, usually store data according with the description of the nested presentation. This presentation makes it then easier to use the databases features.

The nested presentation can be achieved by applying some sort of opposite algorithm to the Flattening Algorithm. This algorithm would nest the objects that are present in the returned graph into a root object. The root object has to be selected by its IRI from the graph, and this identifier has to be one of the supplemented parameters for the construct query.
\todoi{maybe add a more formal name of this algorithm and describe it with a pseudo language, to add more consistency and clearness}
\todoi{add example of the shape we want}

Another possibility how to store the resource is to split it into multiple documents. This would enable us to recognize overlapping data in stored resources of the same type or even across multiple resource types. This would be useful to lower the amount of storage needed and might enable us to more closely emulate a graph database. Splitting the resource into multiple documents would break the biggest advantage of storing data about one resource in one place - fast resource access times. As the objectives state that speed of retrieving the data is more important than the query possibilities its better to store the resource as one document.

\todoi{move to design as this is a bit used db specific}
As document databases use the property names inside the stored JSON, the application of the Compaction algorithm with the right context will shorten the property names.

\section{Service accessibility}
\todoi{add a section about Web APIs and REST }
The service API - what operations we need

\section{Usability possibilities}
\begin{itemize}
	\item How easy it could be to use - analyze steps needed to plug in details caching, querying
	\item Stored resources update
	\item SPARQL to lucene possibilities (from OneNote)
\end{itemize}

\section{Related Work}
\begin{itemize}
	\item Linked Data Fragments - creates local copy of parts of data (fragments) for faster querying (more network local) complex queries get split into easier ones and then evaluated by UNION
\end{itemize}

\section{Optimization - maybe to Conclusion or Implementation}
\begin{itemize}
	\item Offload workload to client
		\subitem Transform of data (required storage implementation details) - i.e. have to rename the properties
		\subitem Analysis of sent queries (renamed properties of documents force a rename in the query) - sparql to lucene in client
\end{itemize}

\chapter{Design}

\section{Web service design}
\section{Query language design}

\chapter{Tests}
\label{chap:Tests}
Describe the test suite by tested queries
Describe test environments - probably my desktop, laptop, vm in azure
\section{Tested queries}
\begin{itemize}
\item create subsections by queries
\end{itemize}

\section{Results}
Compare results from the created service and a sparql enpoint

\section{Test Conclusion}



% Ukázka použití některých konstrukcí LateXu (odkomentujte, chcete-li)
% %%% Ukázka použití některých konstrukcí LaTeXu

\subsection{Ukázka \LaTeX{}u}
\label{ssec:ukazka}

This short subsection serves as an~example of basic \LaTeX{} constructs,
which can be useful for writing a~thesis.

Let us start with lists:

\begin{itemize}
\item The logo of Matfyz is displayed in figure~\ref{fig:mff}.
\item This is subsection~\ref{ssec:ukazka}.
\item Citing literature~\cite{lamport94}.
\end{itemize}

Different kinds of dashes:
red-black (short),
pages 16--22 (middle),
$45-44$ (minus),
and this is --- as you could have expected --- a~sentence-level dash,
which is the longest.
(Note that we have follwed \verb|a| by a~tilde instead of a~space
to avoid line breaks at that place.)

\newtheorem{theorem}{Theorem}
\newtheorem*{define}{Definition}	% Definice nečíslujeme, proto "*"

\begin{define}
A~{\sl Tree} is a connected graph with no cycles.
\end{define}

\begin{theorem}
This theorem is false.
\end{theorem}

\begin{proof}
False theorems do not have proofs.
\end{proof}

\begin{figure}
	\centering
	\includegraphics[width=30mm]{../img/logo.eps}
	\caption{Logo of MFF UK}
	\label{fig:mff}
\end{figure}


\chapter*{Conclusion}
\addcontentsline{toc}{chapter}{Conclusion}


%%% Seznam použité literatury
%%% Seznam použité literatury je zpracován podle platných standardů. Povinnou citační
%%% normou pro diplomovou práci je ISO 690. Jména časopisů lze uvádět zkráceně, ale jen
%%% v kodifikované podobě. Všechny použité zdroje a prameny musí být řádně citovány.

\def\bibname{Bibliography}
\begin{thebibliography}{99}
\addcontentsline{toc}{chapter}{\bibname}

%\bibitem{lamport94}
%  {\sc Lamport,} Leslie.
%  \emph{\LaTeX: A Document Preparation System}.
%  2. vydání.
%  Massachusetts: Addison Wesley, 1994.
%  ISBN 0-201-52983-1.

\end{thebibliography}


%%% Tabulky v diplomové práci, existují-li.
\chapwithtoc{List of Tables}

%%% Použité zkratky v diplomové práci, existují-li, včetně jejich vysvětlení.
\chapwithtoc{List of Abbreviations}

%%% Přílohy k diplomové práci, existují-li (různé dodatky jako výpisy programů,
%%% diagramy apod.). Každá příloha musí být alespoň jednou odkazována z vlastního
%%% textu práce. Přílohy se číslují.
\chapwithtoc{Attachments}

\openright
\end{document}
